\documentclass[11pt]{article}
\usepackage[margin=1in]{geometry}
\usepackage{amsmath,amssymb,bm}
\usepackage{booktabs}
\usepackage{array}
\usepackage{xcolor}
\usepackage{enumitem}
\usepackage{algorithm}
\usepackage{algpseudocode}
\usepackage{hyperref}

\title{Brief Mathematical Workflow for DAD, Myopic Entropy Reduction, and Baselines}
\author{}
\date{}

\newcommand{\thetaVec}{\bm{\theta}}
\newcommand{\Mvec}{\bm{M}}
\newcommand{\Kvec}{\bm{K}}
\newcommand{\xivec}{\xi}
\newcommand{\hist}{h}
\newcommand{\Npart}{N}
\newcommand{\Normal}{\mathcal{N}}
\newcommand{\Aset}{\mathcal{A}}

\begin{document}
\maketitle

\section{Problem Setting}

The unknown parameter is the bus-level inertia and damping vector
\begin{equation}
\thetaVec=(\Mvec,\Kvec)\in \mathbb{R}^{28},
\qquad
\Mvec=(M_1,\ldots,M_{14}),\quad
\Kvec=(K_1,\ldots,K_{14}).
\end{equation}

At each experimental step $t=1,\ldots,T$, the method selects one design
\begin{equation}
\xi_t=(a_t,b_t,d_t),
\end{equation}
where $a_t$ is the probing amplitude, $b_t$ is the selected bus location, and $d_t$ is the probing duration. In the current setting, $d_t=0.2$ s is fixed.

The history before step $t$ is
\begin{equation}
 h_{t-1}=\{(\xi_1,y_{\mathrm{obs},1}),\ldots,(\xi_{t-1},y_{\mathrm{obs},t-1})\}.
\end{equation}

The offline train bank contains prior particles
\begin{equation}
\{\thetaVec_n\}_{n=1}^{N},
\end{equation}
and precomputed clean responses for feasible design prefixes:
\begin{equation}
 y_{\mathrm{clean}}(\thetaVec_n,\xi_{1:t}).
\end{equation}
The noisy observation model is
\begin{equation}
 y_{\mathrm{obs},t}
 = y_{\mathrm{clean}}(\thetaVec,\xi_{1:t})+\epsilon_t,
 \qquad
 \epsilon_t\sim \Normal(0,\sigma_y^2).
\end{equation}
Thus, for a particle $\thetaVec_n$,
\begin{equation}
 p(y_{\mathrm{obs},t}\mid \thetaVec_n,\xi_{1:t})
 =
 \Normal\!\big(y_{\mathrm{obs},t};
 y_{\mathrm{clean}}(\thetaVec_n,\xi_{1:t}),
 \sigma_y^2\big).
\end{equation}

\section{Mathematical Basics}

\subsection{Entropy Formula}

Using the train particles as finite posterior support, the belief after history $h_t$ is represented by weights
\begin{equation}
 p(\thetaVec\mid h_t)\approx \sum_{n=1}^{N} w_n^{(t)}\delta_{\thetaVec_n}(\thetaVec),
 \qquad
 \sum_{n=1}^{N}w_n^{(t)}=1.
\end{equation}

The particle entropy is
\begin{equation}
 H(w^{(t)})
 =
 -\sum_{n=1}^{N}w_n^{(t)}\log w_n^{(t)}.
\end{equation}

With a uniform prior over train particles,
\begin{equation}
 w_n^{(0)}=\frac{1}{N},
 \qquad
 H(w^{(0)})=\log N.
\end{equation}

After observing a history $h_T$, the realized entropy reduction is
\begin{equation}
 \Delta H(h_T)
 =
 H(w^{(0)})-H(w^{(T)}).
\end{equation}

\subsection{Information Metrics in Evaluation}

Entropy reduction $\Delta H(h_T)$ is computed after one realized history is observed. Expected information gain (EIG) averages this quantity over possible future histories:
\begin{equation}
 \mathrm{EIG}(\pi)
 =
 \mathbb{E}_{h_T\sim \pi}
 \left[
 H(w^{(0)})-H(w^{(T)})
 \right]
 =
 \mathbb{E}_{h_T\sim \pi}[\Delta H(h_T)].
\end{equation}

In finite test evaluation, we report the sample mean realized entropy reduction
\begin{equation}
 \widehat{\Delta H}
 =
 \frac{1}{N_{\mathrm{test}}}
 \sum_{m=1}^{N_{\mathrm{test}}}
 \Delta H(h_T^{(m)}).
\end{equation}
This statistic is a realized-information diagnostic on test rollouts, not an exact estimator of policy EIG.

\subsection{sPCE Formula (Primary Information Score)}

The sequential prior contrastive estimation (sPCE) objective is a contrastive lower-bound-style estimator of EIG. For one rollout generated from a true particle $\thetaVec_0$ and $L$ contrastive particles $\thetaVec_1,\ldots,\thetaVec_L$, define
\begin{equation}
 R_{\mathrm{sPCE}}(h_T)
 =
 \log p(h_T\mid \thetaVec_0)
 -
 \log\left(
 \frac{1}{L+1}
 \sum_{\ell=0}^{L}p(h_T\mid \thetaVec_\ell)
 \right).
\end{equation}
In this workflow, sPCE is the primary information metric for method comparison; $\Delta H$ is reported as a complementary posterior-uncertainty diagnostic.

Because the design choices are already contained in the rollout history, the likelihood term can be computed from the observation likelihood:
\begin{equation}
 \log p(h_T\mid \thetaVec_\ell)
 =
 \sum_{t=1}^{T}
 \log
 \Normal\!\left(
 y_{\mathrm{obs},t};
 y_{\mathrm{clean}}(\thetaVec_\ell,\xi_{1:t}),
 \sigma_y^2
 \right).
\end{equation}

\section{DAD with Different Reward Functions}

DAD uses a trainable adaptive policy
\begin{equation}
 \xi_t\sim \pi_\phi(\cdot\mid h_{t-1}),
\end{equation}
where $\phi$ denotes neural-network parameters. DAD has an offline training stage and a fast evaluation stage.

\subsection{DAD-sPCE}

DAD-sPCE trains the policy using the sPCE reward:
\begin{equation}
 R(h_T)=R_{\mathrm{sPCE}}(h_T).
\end{equation}
The policy objective is
\begin{equation}
 \max_\phi\; \mathbb{E}_{h_T\sim \pi_\phi}[R_{\mathrm{sPCE}}(h_T)].
\end{equation}
Equivalently, the loss can be written as the negative reward objective. With REINFORCE, one common loss form is
\begin{equation}
 \mathcal{L}_{\mathrm{DAD\mbox{-}sPCE}}(\phi)
 =
 -\big(R_{\mathrm{sPCE}}(h_T)-b\big)
 \sum_{t=1}^{T}
 \log \pi_\phi(\xi_t\mid h_{t-1}),
\end{equation}
where $b$ is a baseline used to reduce gradient variance.

\paragraph{Training mode used in practice.}
The DAD training mode is REINFORCE with sPCE reward.

\subsection{DAD-$\Delta H$}

DAD-$\Delta H$ trains the policy using direct posterior entropy reduction:
\begin{equation}
 R(h_T)=R_{\Delta H}(h_T)=H(w^{(0)})-H(w^{(T)}).
\end{equation}
The objective is
\begin{equation}
 \max_\phi\;\mathbb{E}_{h_T\sim \pi_\phi}[\Delta H(h_T)].
\end{equation}
Using REINFORCE, the loss can be written as
\begin{equation}
 \mathcal{L}_{\mathrm{DAD}\mbox{-}\Delta H}(\phi)
 =
 -\big(R_{\Delta H}(h_T)-b\big)
 \sum_{t=1}^{T}
 \log \pi_\phi(\xi_t\mid h_{t-1}).
\end{equation}

\subsection{DAD Pseudocode}

\begin{algorithm}[H]
\caption{DAD Training with sPCE or $\Delta H$ Reward}
\begin{algorithmic}[1]
\State Input train bank, policy $\pi_\phi$, reward type $R\in\{R_{\mathrm{sPCE}},R_{\Delta H}\}$
\For{each training iteration}
    \State Sample a true train particle $\thetaVec_0$ and initialize $h_0=\emptyset$
    \For{$t=1,\ldots,T$}
        \State Sample/select design $\xi_t\sim\pi_\phi(\cdot\mid h_{t-1})$
        \State Read/generate noisy observation $y_{\mathrm{obs},t}$ from train bank
        \State Update history $h_t=h_{t-1}\cup\{(\xi_t,y_{\mathrm{obs},t})\}$
    \EndFor
    \State Compute trajectory reward $R(h_T)$
    \State Update $\phi$ by minimizing $-(R(h_T)-b)\sum_t \log\pi_\phi(\xi_t\mid h_{t-1})$
\EndFor
\State Output trained policy $\pi_{\phi^*}$
\end{algorithmic}
\end{algorithm}

At evaluation time, the trained policy is frozen:
\begin{equation}
 \xi_t=\pi_{\phi^*}(h_{t-1}).
\end{equation}
The final performance is evaluated using the same metrics as the baselines: sPCE score, $\Delta H$, and MSE of $\thetaVec$.

\section{Myopic Entropy-Reduction Baseline}

The myopic baseline is adaptive but not trained. It has no trainable parameters and no training loss. Instead, it performs online Bayesian updating and greedy one-step design scoring during evaluation.

\subsection{Posterior Weight Update}

At the start,
\begin{equation}
 w_n^{(0)}=\frac{1}{N}.
\end{equation}
After observing $y_{\mathrm{obs},t}$ under selected prefix $\xi_{1:t}$, update particle weights by
\begin{equation}
 \tilde w_n^{(t)}
 =
 w_n^{(t-1)}
 \Normal\!\left(
 y_{\mathrm{obs},t};
 y_{\mathrm{clean}}(\thetaVec_n,\xi_{1:t}),
 \sigma_y^2
 \right),
\end{equation}
\begin{equation}
 w_n^{(t)}
 =
 \frac{\tilde w_n^{(t)}}{\sum_{r=1}^{N}\tilde w_r^{(t)}}.
\end{equation}

\subsection{Candidate Design Score}

Before choosing the next design at step $t$, the future observation is unknown. For each feasible candidate design $\xi\in\Aset_t$, define the candidate prefix
\begin{equation}
 \xi_{1:t}^{\mathrm{cand}}=(\xi_1,\ldots,\xi_{t-1},\xi).
\end{equation}

For each particle, read the clean response
\begin{equation}
 \mu_n(\xi)
 =
 y_{\mathrm{clean}}(\thetaVec_n,\xi_{1:t}^{\mathrm{cand}}).
\end{equation}

A simple predicted observation is the posterior predictive mean
\begin{equation}
 \hat y(\xi)
 =
 \sum_{n=1}^{N}w_n^{(t-1)}\mu_n(\xi).
\end{equation}

The hypothetical updated weight under candidate $\xi$ is
\begin{equation}
 \bar w_n(\xi)
 \propto
 w_n^{(t-1)}
 \Normal\!\left(
 \hat y(\xi);
 \mu_n(\xi),
 \sigma_y^2
 \right),
\end{equation}
with normalization $\sum_n \bar w_n(\xi)=1$.

The one-step estimated entropy-reduction score is
\begin{equation}
 \widehat{\Delta H}_t(\xi\mid h_{t-1})
 =
 H(w^{(t-1)})-H(\bar w(\xi)).
\end{equation}

The selected design is
\begin{equation}
 \xi_t^*
 =
 \arg\max_{\xi\in\Aset_t}
 \widehat{\Delta H}_t(\xi\mid h_{t-1}).
\end{equation}

After $\xi_t^*$ is selected, the method observes the real test response $y_{\mathrm{obs},t}$ and updates $w^{(t)}$ using the posterior weight update above.

\subsection{Myopic Pseudocode}

\begin{algorithm}[H]
\caption{Myopic Entropy-Reduction Baseline}
\begin{algorithmic}[1]
\State Input train particles, train clean-response table, test noisy-response table
\State Initialize $w_n^{(0)}=1/N$ and $h_0=\emptyset$
\For{$t=1,\ldots,T$}
    \For{each feasible candidate design $\xi\in\Aset_t$}
        \State Compute $\mu_n(\xi)=y_{\mathrm{clean}}(\thetaVec_n,\xi_{1:t}^{\mathrm{cand}})$ for all train particles
        \State Compute predicted observation $\hat y(\xi)=\sum_n w_n^{(t-1)}\mu_n(\xi)$
        \State Compute hypothetical weights $\bar w(\xi)$
        \State Compute score $\widehat{\Delta H}_t(\xi\mid h_{t-1})=H(w^{(t-1)})-H(\bar w(\xi))$
    \EndFor
    \State Select $\xi_t^*=\arg\max_{\xi}\widehat{\Delta H}_t(\xi\mid h_{t-1})$
    \State Observe real test response $y_{\mathrm{obs},t}$ for selected design
    \State Update posterior weights $w^{(t)}$ using the real observation
    \State Update history $h_t=h_{t-1}\cup\{(\xi_t^*,y_{\mathrm{obs},t})\}$
\EndFor
\State Output selected design sequence and final $\Delta H=H(w^{(0)})-H(w^{(T)})$
\end{algorithmic}
\end{algorithm}

\section{Other Baselines}

\subsection{Random Baseline}

The random baseline selects a valid design randomly at each step:
\begin{equation}
 \xi_t\sim \mathrm{Uniform}(\Aset_t).
\end{equation}
It has no training and no entropy-based decision score. After the full sequence is observed, the same posterior update is applied and the final $\Delta H$ is computed for evaluation.

\subsection{Fixed-Design Baseline}

The fixed-design baseline uses a pre-specified sequence
\begin{equation}
 \xi_{1:T}^{\mathrm{fixed}}
 =
 (\xi_1^{\mathrm{fixed}},\ldots,\xi_T^{\mathrm{fixed}}).
\end{equation}
In this document, a recommended fixed sequence is an equal-interval bus design (e.g., for $T=3$, buses spread across the network with fixed amplitude and duration) so that the baseline is deterministic and physically interpretable. Any deterministic pre-specified sequence can be substituted in implementation.
It is non-adaptive because
\begin{equation}
 \xi_t^{\mathrm{fixed}}
 \text{ does not depend on } h_{t-1}.
\end{equation}
After observing the responses from this fixed sequence, posterior weights are updated and $\Delta H$ is computed.

\subsection{Baseline Pseudocode}

\begin{algorithm}[H]
\caption{Fixed Baseline Evaluation}
\begin{algorithmic}[1]
\State Choose a deterministic fixed sequence $\xi_{1:T}^{\mathrm{fixed}}$ (equal-interval bus rule)
\State For each test system, read $y_{\mathrm{obs},1:T}$ under $\xi_{1:T}^{\mathrm{fixed}}$
\State Compute posterior weights on train particles using Gaussian updates with $y_{\mathrm{clean}}$
\State Report sPCE, $\Delta H$, and MSE with the shared evaluation routine
\end{algorithmic}
\end{algorithm}

\section{Summary Table}

\begin{table}[h]
\centering
\renewcommand{\arraystretch}{1.25}
\begin{tabular}{p{3.2cm}ccp{6.6cm}}
\toprule
Method & Adaptive? & Training? & Objective / score \\
\midrule
Fixed design & No & No & Uses a deterministic equal-interval sequence; evaluated by shared scores (sPCE, $\Delta H$, MSE). \\
Random & No & No & Random valid design selection; mainly a weak sanity baseline. \\
Myopic entropy-reduction baseline & Yes & No & Greedy online score $\widehat{\Delta H}_t(\xi\mid h_{t-1})$ for each candidate design. \\
DAD-sPCE & Yes & Yes & Offline policy training with sPCE reward $R_{\mathrm{sPCE}}$ and policy-gradient loss. \\
DAD-$\Delta H$ & Yes & Yes & Offline policy training with direct entropy-reduction reward $R_{\Delta H}=H(w^{(0)})-H(w^{(T)})$. \\
\bottomrule
\end{tabular}
\end{table}

\paragraph{Key distinction.}
Myopic uses entropy reduction as an online design score during evaluation. DAD uses sPCE or $\Delta H$ as an offline training reward to learn a policy. Fixed and random baselines do not optimize an information objective when selecting designs; they are evaluated with the same shared metrics (sPCE, $\Delta H$, MSE).

\end{document}
